% =============================================================================
% Ready-to-paste LaTeX: repeated cross-validation, confidence intervals and both
% significance tests for all headline comparisons.
%
% Numbers from experiments/evaluation/exp_cv_corrected.py (results/cv_corrected.json):
% 10 x 5 repeated GroupKFold, nested-CV-tuned C, macro-F1 on the out-of-fold predictions
% of each repeat pooled together, seed 42. The effect of tuning C (Hyper-parameter
% selection paragraph) is from results/statistical_power.json, default vs tuned.
%
% Do NOT reintroduce the older per-fold figures (0.397 / 0.394 / 0.350 / 0.345 / 0.343 /
% 0.335; 8 genres 0.300 / 0.257). They are biased low by 0.02-0.03 because rare genres
% are absent from many single test folds, and they are superseded here (progress log
% section 21). statistical_power.json keeps them as the record of that protocol.
%
% Bib: nadeau2003inference (in bib/library.bib).
%
% Placement: Methods (the "Evaluation protocol" paragraph) + Results (the tables and the
% discussion), or collect the whole thing as a "Statistical reliability" subsection at the
% end of Results.
% =============================================================================

% ---------------------------------------------------------------- METHODS ----
\paragraph{Evaluation protocol and statistical reliability.}
All genre results are obtained under cross-validation grouped by film, so that clips
from the same soundtrack never appear in both the training and the test partition; this
is essential, since an ungrouped split lets a classifier recognise the film rather than
the genre (Section~\ref{subsec:leakage}). A single five-fold run, however, yields only
five score samples, which is too few to distinguish a difference of a few macro-$F_1$
points from the variation caused by the particular partition. The evaluation is
therefore repeated: the assignment of films to folds is re-randomised ten times, giving
$10\times5=50$ leakage-free folds, and every method is scored on exactly the same folds
so that all comparisons are paired.

Grouping by film has a consequence for how the score is computed. A genre represented by
only a few films is absent from many test folds entirely -- Comedy, with four films, has
no clip in 17 of the 50, Horror in 15 -- and a macro-$F_1$ computed fold by fold scores
such a genre as zero there, whatever the model predicts. Scores are therefore computed
once per repeat: the five test folds of a repeat together contain every clip exactly
once, so their out-of-fold predictions are pooled and scored together. The per-fold
alternative lowers the score of every model by about $0.02$--$0.03$ while leaving the
ranking of the methods unchanged. Where the genre classifier is trained on \emph{predicted} emotions, the
emotions of the training clips are themselves predicted out-of-fold, by an inner
cross-validation over the training films, so that the classifier is trained on the same
kind of input it will be tested on.

Uncertainty on a single method is given as a $95\,\%$ interval from a bootstrap over
films: the 41 films are resampled with replacement 2000 times and the pooled predictions
re-scored on each resample. The films, not the clips, are resampled, because clips of the
same film are not independent. Two methods are compared by two tests, which answer
slightly different questions. The \emph{paired film bootstrap} applies the same resampling
to the difference between two methods and asks how much it depends on which films happen
to make up the corpus, with the trained models held fixed. The Nadeau--Bengio corrected
resampled $t$-test \citep{nadeau2003inference}, applied to the fifty per-fold scores, also
accounts for which films the models happened to be trained on, and is the more conservative
of the two. An ordinary paired $t$-test would be inadmissible here: it treats the fifty
fold-differences as independent although their training sets overlap almost completely.
The correction inflates the variance estimate by the ratio of test-set to training-set
size,
\[
  t \;=\; \frac{\bar{d}}{\sqrt{\left(\dfrac{1}{n}+\dfrac{n_{\text{test}}}{n_{\text{train}}}\right)\sigma^2_{d}}},
\]
where $\bar{d}$ and $\sigma^2_d$ are the mean and variance of the per-fold differences
and $n=50$. Because only the $1/n$ term shrinks when repeats are added, this test has a
lower bound on the attainable $p$-value for a given corpus; that bound is reported as
$p_{\lim}$ and is what makes it possible to distinguish a comparison that needs more
resampling from one that needs more data. Where the two tests disagree, both are reported.

\paragraph{Hyper-parameter selection.}
The inverse regularisation strength $C$ of the logistic regression is not left at
scikit-learn's default but selected per training fold by an inner cross-validation, again
grouped by film -- a nested cross-validation. This matters here: with roughly $330$ clips
against embedding dimensionalities of $128$ to $768$, the default $C=1$ under-regularises
every feature set, and it penalises the high-dimensional embeddings most (tuning raises
the AST baseline by $0.025$ and the CLAP baseline by $0.033$ macro-$F_1$, against $0.018$
for the eleven emotion features). Selecting $C$ on the test folds would inflate the
reported scores, and leaving it untuned would understate the baselines; the nested
procedure avoids both. The most frequently selected value is $C=0.003$ for every feature
set except the eight-dimensional PCA control, for which it is $0.1$.

% ---------------------------------------------------------------- RESULTS ----
\subsection{Statistical Reliability of the Comparisons}
\label{subsec:statistical_power}

Table~\ref{tab:power_arms} reports each approach on the five-genre subset under the
repeated protocol, and Table~\ref{tab:power_cmp} the corresponding paired comparisons.

\begin{table}[t]
  \centering
  \caption{Genre classification on the five-genre subset ($n=329$ clips, 41 films)
    under $10\times5$ repeated GroupKFold with nested-CV-tuned $C$. Macro-$F_1$ on the
    pooled out-of-fold predictions, averaged over the ten repeats, with a $95\,\%$
    interval from a bootstrap over films.}
  \label{tab:power_arms}
  \begin{tabular}{lcc}
    \hline
    Approach & Macro-$F_1$ & $95\,\%$ CI \\
    \hline
    Ground-truth emotion (11) $\rightarrow$ genre (ceiling)         & $0.428$ & $[0.371,\,0.476]$ \\
    VGGish $\rightarrow$ predicted emotion (11) $\rightarrow$ genre & $0.417$ & $[0.365,\,0.460]$ \\
    VGGish-128 $\rightarrow$ genre (direct)                         & $0.377$ & $[0.325,\,0.414]$ \\
    PCA-8(VGGish) $\rightarrow$ genre (control)                     & $0.372$ & $[0.316,\,0.416]$ \\
    AST-768 $\rightarrow$ genre                                     & $0.371$ & $[0.326,\,0.398]$ \\
    MIR-103 (hand-crafted) $\rightarrow$ genre                      & $0.363$ & $[0.318,\,0.398]$ \\
    MusiCNN-200 $\rightarrow$ genre                                 & $0.363$ & $[0.311,\,0.401]$ \\
    CLAP-512 $\rightarrow$ genre                                    & $0.361$ & $[0.313,\,0.400]$ \\
    Most-frequent baseline                                          & $0.168$ & $[0.146,\,0.185]$ \\
    \hline
  \end{tabular}
\end{table}

\begin{table}[t]
  \centering
  \caption{Paired comparisons on the five-genre subset. $\Delta$ is the macro-$F_1$
    difference on the pooled predictions with a $95\,\%$ interval from the film bootstrap;
    $p_{\text{boot}}$ is the paired film bootstrap, $p_{\text{NB}}$ the Nadeau--Bengio
    corrected resampled $t$-test on the per-fold scores, and $p_{\lim}$ the smallest
    $p_{\text{NB}}$ this corpus allows with an unlimited number of repeats; ``won'' is the
    number of the ten repeats in which the first method scores higher.}
  \label{tab:power_cmp}
  \begin{tabular}{lcccccc}
    \hline
    Comparison & $\Delta$ & $95\,\%$ CI & $p_{\text{boot}}$ & $p_{\text{NB}}$ & $p_{\lim}$ & won \\
    \hline
    Predicted emotion vs.\ VGGish direct & $+0.040$ & $[+0.001,\,+0.084]$ & $\mathbf{0.042}$ & $0.151$ & $0.130$ & 10 \\
    Predicted emotion vs.\ AST-768       & $+0.047$ & $[+0.004,\,+0.094]$ & $\mathbf{0.030}$ & $0.110$ & $0.091$ & 10 \\
    Predicted emotion vs.\ CLAP-512      & $+0.055$ & $[+0.010,\,+0.103]$ & $\mathbf{0.017}$ & $0.107$ & $0.088$ & 10 \\
    Predicted emotion vs.\ MusiCNN-200   & $+0.054$ & $[+0.008,\,+0.106]$ & $\mathbf{0.018}$ & $0.073$ & $0.057$ & 10 \\
    Predicted emotion vs.\ MIR-103       & $+0.054$ & $[+0.003,\,+0.105]$ & $\mathbf{0.036}$ & $0.070$ & $0.055$ & 10 \\
    Predicted emotion vs.\ PCA-8 control & $+0.045$ & $[+0.006,\,+0.087]$ & $\mathbf{0.029}$ & $0.122$ & $0.102$ & 10 \\
    Predicted emotion vs.\ ground truth  & $-0.011$ & $[-0.035,\,+0.012]$ & $0.325$ & $0.497$ & $0.477$ & 3 \\
    Ground-truth emotion vs.\ CLAP-512   & $+0.066$ & $[+0.027,\,+0.108]$ & $\mathbf{0.001}$ & $\mathbf{0.036}$ & $0.025$ & 10 \\
    PCA-8 control vs.\ VGGish direct     & $-0.005$ & $[-0.034,\,+0.024]$ & $0.771$ & $0.804$ & $0.795$ & 3 \\
    AST-768 vs.\ CLAP-512                & $+0.008$ & $[-0.030,\,+0.045]$ & $0.663$ & $0.752$ & $0.741$ & 7 \\
    \hline
  \end{tabular}
\end{table}

Four findings follow. First, the ordering of the approaches is \emph{stable}: the
predicted-emotion pipeline scores higher than every direct representation, and higher
than the PCA-8 control, in all ten repeats, by $0.040$ to $0.055$, and the bootstrap
interval of every one of these differences excludes zero. Second, the direct
representations are statistically indistinguishable from one another -- $0.361$ to
$0.377$, with $p=0.66$ between AST and CLAP -- so the comparison is not against one
arbitrarily chosen or poorly suited audio representation: the emotion features outperform
all five, spanning general-audio, audio--text, music-specific and hand-crafted
features. Compression on its own does not account for the advantage, since the PCA-8
control is no better than the embedding it compresses ($\Delta=-0.005$, $p=0.77$).
Third, predicted and ground-truth emotions are themselves indistinguishable
($\Delta=-0.011$, $p=0.33$), meaning the emotion values estimated from audio are as
informative for genre as the human ratings -- the central prerequisite for the two-stage
pipeline to be usable without manually annotated emotion data.

Fourth, whether the advantage is \emph{significant} depends on which question is asked.
Under the film bootstrap every comparison of the pipeline against a direct representation
or the control is significant ($p$ between $0.017$ and $0.042$); under the more
conservative Nadeau--Bengio test they remain borderline ($p$ between $0.07$ and $0.15$).
The effect sizes are the same under both; what differs is that the second test also
charges for the variation that comes from which films were used for training. The
advantage is therefore reported as consistent and substantial, significant under one test
and borderline under the stricter one, rather than as an unqualified significant result.

That remaining uncertainty is a property of the corpus rather than of the evaluation
procedure. Because only the $1/n$ term of the corrected test shrinks with additional
repeats, the attainable $p$-value has a floor for a given dataset; at ten repeats the
test is already within about $0.02$ of it, and for every comparison of the pipeline
against a single representation or the control that floor lies above $0.05$
($p_{\lim}$ between $0.055$ and $0.130$). No number of repeats could make these
comparisons significant under this test; only more films could. This agrees with the
learning-curve analysis in Section~\ref{subsec:learning_curve}, where neither model had
reached a plateau at the full corpus size.

On the full eight-genre task the advantage is significant under both tests. Ground-truth
emotion features attain $0.321$ $[0.280,\,0.352]$ against $0.276$ $[0.243,\,0.296]$ for
the AST embedding, a difference of $+0.046$ ($p_{\text{boot}}=0.008$,
$p_{\text{NB}}=0.031$). The full pipeline, with emotions predicted from VGGish, reaches
$0.318$, ahead of AST by $+0.043$ ($p_{\text{boot}}=0.012$, $p_{\text{NB}}=0.049$) and of
the VGGish embedding it starts from ($0.268$) by $+0.051$ ($p_{\text{boot}}<0.001$,
$p_{\text{NB}}=0.011$); each of these differences holds in all ten repeats, against a
most-frequent baseline of $0.102$. Under the earlier single five-fold protocol the
ground-truth comparison with AST produced $p=0.62$, which illustrates how little a
five-fold run can say about a difference of this size.
